\documentclass[11pt]{article}

\usepackage[margin=1in]{geometry}
\usepackage{amsmath,amssymb,amsfonts,mathtools,amsthm}
\usepackage{bm}
\usepackage{booktabs,array,longtable,tabularx}
\usepackage{enumitem}
\usepackage{xcolor}
\usepackage{xurl}
\usepackage[hidelinks]{hyperref}

\numberwithin{equation}{section}
\setlist{nosep,leftmargin=*}
\setlength{\emergencystretch}{3em}

\newcommand{\CP}{\mathbb{CP}}
\newcommand{\CC}{\mathbb C}
\newcommand{\QQ}{\mathbb Q}
\newcommand{\cO}{\mathcal O}
\newcommand{\sltwo}{\mathfrak{sl}_2(\mathbb C)}
\newcommand{\Ktr}{K_{\mathrm{tr}}}
\newcommand{\dd}{\mathrm d}
\newcommand{\Tr}{\operatorname{Tr}}
\newcommand{\Sym}{\operatorname{Sym}}
\newcommand{\Span}{\operatorname{Span}}
\newcommand{\diag}{\operatorname{diag}}
\newcommand{\rank}{\operatorname{rank}}
\newcommand{\ad}{\operatorname{ad}}
\newcommand{\Div}{\operatorname{Div}}
\newcommand{\status}[1]{\textsf{#1}}
\newcommand{\nopromote}{\texttt{physics\_promotion\_allowed=false}}

\definecolor{darkgreen}{RGB}{0,105,55}
\definecolor{darkred}{RGB}{150,25,25}
\definecolor{darkblue}{RGB}{20,65,140}

\newtheorem{theorem}{Theorem}[section]
\newtheorem{proposition}[theorem]{Proposition}
\newtheorem{lemma}[theorem]{Lemma}
\newtheorem{corollary}[theorem]{Corollary}
\newtheorem{definition}[theorem]{Definition}
\newtheorem{remark}[theorem]{Remark}
\newtheorem{gate}[theorem]{Promotion gate}

\title{AP-E2: Exact Discriminant--Contact Regression\\
First-Principles Projective Charts, Killing Normalization,\\
and the Fail-Closed Route-E Convention Boundary}
\author{Route F research note}
\date{14 July 2026}

\begin{document}
\maketitle

\begin{abstract}
This note completes AP-E2 as an exact, non-promoting regression theorem.  On
the conditional carrier
\(H^0(\CP^1,T_{\CP^1})\simeq H^0(\CP^1,\cO(2))\simeq\sltwo\), let
\(p(u,v)=au^2+buv+cv^2\) and
\(
 A_p=\left(\begin{smallmatrix}-b/2&-c\\a&b/2\end{smallmatrix}\right)
\).
The projective convention is fixed, rather than guessed, by the identity
\(\det(z,A_pz)=p(z)\).  Direct multiplication and the adjoint-trace
normalization then give
\[
 A_p^2=\frac{\Delta(p)}4I,
 \qquad B_{\rm Kill}(A_p,A_p)=2\Delta(p),
 \qquad (p,p)_2=-2\Delta(p),
 \]
where \(\Delta=b^2-4ac\).  Thus, for a nonzero section, two distinct
projective zeros, a regular semisimple generator, nonzero Killing norm, and
nonzero Route-E contact are equivalent.  The boundary \(\Delta=0\) is the
nonzero nilpotent double-zero cone; the zero section is a separate excluded
stratum.

The note also resolves the convention-sensitive factor of two.  In Route-E
spherical coordinates \(\bm x=(x_+,x_0,x_-)^T\),
\[
 B=2\sqrt3\,\bm x^T\Ktr\bm x,
 \]
whereas in normalized polynomial coefficients
\(\bm y=(a,b/\sqrt2,c)^T\),
\[
 B=4\sqrt3\,\bm y^T\Ktr\bm y.
\]
These formulas are compatible because \(\bm y=R\bm x/\sqrt2\).  A
deterministic 30-check audit tests exact rational identities, projective
charts, \(SL(2)\) covariance, the degenerate boundary, Route-E source hashes,
100-digit complex samples, and four deliberately wrong conventions.  All
checks pass, but \nopromote: AP-E2 preserves an
algebraic theorem and does not derive H3$^+$, a physical two-center
interpretation, or the AP-E3 level \(k=2\).
\end{abstract}

\tableofcontents

\section{Scope, assumptions, and final logical status}

AP-E2 has one narrow purpose: freeze the exact discriminant/contact theorem
against later changes of chart, basis, or normalization.  It does not add a
new action and does not infer dynamics from an invariant polynomial.

The only geometric input is the conditional Route-E branch
\begin{equation}
 \Sigma\simeq\CP^1,
 \qquad T_{\CP^1}\simeq\cO(2),
 \qquad V:=H^0(\CP^1,T_{\CP^1})\simeq\sltwo.
 \label{eq:conditional-carrier}
\end{equation}
The word ``conditional'' is essential: the algebra below is unconditional
after Eq.~\eqref{eq:conditional-carrier} is assumed, but AP-E2 does not derive
the microscopic origin of that carrier or of the stronger contact hypothesis
H3$^+$.

We distinguish four statuses.
\begin{longtable}{>{\raggedright\arraybackslash}p{0.20\textwidth}
                  >{\raggedright\arraybackslash}p{0.71\textwidth}}
\toprule
Status & Meaning \\
\midrule
\endfirsthead
\toprule
Status & Meaning \\
\midrule
\endhead
\status{exact theorem} & An identity proved from the displayed definitions
over \(\CC\), with every normalization fixed.\\
\status{regression pass} & A deterministic machine check of a theorem,
boundary case, source hash, or known convention failure.\\
\status{conditional carrier} & A conclusion valid on
Eq.~\eqref{eq:conditional-carrier}; it is not a derivation of that equation.\\
\status{not promoted} & A geometric resemblance has not been turned into a
physical state-space, action, spectrum, or family-selection mechanism.\\
\bottomrule
\end{longtable}

\begin{gate}[AP-E2 non-promotion]
A green AP-E2 regression may certify the exact invariant identity and its
regular/nilpotent stratification.  It may not identify the two zeros with two
particles, opposite energies, or two families; it may not promote H3 to
H3$^+$; and it may not supply the level-two prequantum bundle required by
AP-E3.  The output Boolean must therefore remain
\nopromote.
\end{gate}

\section{First principles on \texorpdfstring{$\CP^1$}{CP1}}

\subsection{Two charts force a quadratic vector field}

Let \([u:v]\) be homogeneous coordinates on \(\CP^1\).  On
\(U_u=\{u\ne0\}\), set
\begin{equation}
 \xi=\frac vu,
 \label{eq:xi-chart}
\end{equation}
and write a holomorphic vector field as
\begin{equation}
 X=q(\xi)\partial_\xi.
 \label{eq:q-vector}
\end{equation}
On \(U_v=\{v\ne0\}\), set \(\eta=u/v=1/\xi\).  Since
\(\partial_\xi=-\eta^2\partial_\eta\), the same field is
\begin{equation}
 X=-\eta^2q(1/\eta)\partial_\eta.
 \label{eq:eta-transform}
\end{equation}
Regularity at \(\eta=0\) forbids powers of \(\xi\) above two.  Therefore
\begin{equation}
 q(\xi)=a+b\xi+c\xi^2,
 \qquad a,b,c\in\CC.
 \label{eq:q-quadratic}
\end{equation}
Conversely, Eq.~\eqref{eq:eta-transform} becomes
\begin{equation}
 X=-(a\eta^2+b\eta+c)\partial_\eta,
 \label{eq:eta-polynomial}
\end{equation}
which is regular.  Hence Eq.~\eqref{eq:q-quadratic} lists every global
holomorphic vector field and proves directly that the section space has
complex dimension three.

The basis
\begin{equation}
 E=\partial_\xi,
 \qquad H=-2\xi\partial_\xi,
 \qquad F=-\xi^2\partial_\xi
 \label{eq:sl2-vector-basis}
\end{equation}
obeys
\begin{equation}
 [H,E]=2E,
 \qquad [H,F]=-2F,
 \qquad [E,F]=H.
 \label{eq:sl2-brackets}
\end{equation}
This constructs the isomorphism
\(H^0(\CP^1,T_{\CP^1})\simeq\sltwo\) without dimension matching alone.

\subsection{Homogenization and the projective matrix}

Homogenizing Eq.~\eqref{eq:q-quadratic} gives
\begin{equation}
 p(u,v)=au^2+buv+cv^2,
 \qquad p(1,\xi)=q(\xi).
 \label{eq:binary-quadratic}
\end{equation}
For \(z=(u,v)^T\), define
\begin{equation}
 A_p=
 \begin{pmatrix}
 -b/2&-c\\
 a&b/2
 \end{pmatrix}
 \in\sltwo.
 \label{eq:projective-matrix}
\end{equation}
This sign convention is fixed by the following identity.

\begin{lemma}[Section--eigenline identity]
For every \(z=(u,v)^T\),
\begin{equation}
 \boxed{\det(z,A_pz)=p(u,v).}
 \label{eq:wedge-identity}
\end{equation}
Consequently, \([u:v]\) is a zero of \(p\) if and only if it is an
eigenline of \(A_p\).
\end{lemma}

\begin{proof}
Equation~\eqref{eq:projective-matrix} gives
\begin{equation}
 A_pz=
 \begin{pmatrix}-bu/2-cv\\au+bv/2\end{pmatrix}.
\end{equation}
Taking the determinant of the matrix with columns \(z,A_pz\) gives
\begin{align}
 \det(z,A_pz)
 &=u(au+bv/2)-v(-bu/2-cv)\\
 &=au^2+buv+cv^2=p(u,v).
\end{align}
The determinant vanishes precisely when \(z\) and \(A_pz\) are linearly
dependent, which is the projective eigenline condition.
\end{proof}

\begin{remark}[Why the chart matters]
On \(u\ne0\), Eq.~\eqref{eq:wedge-identity} reads
\(u^2q(\xi)=0\).  At infinity, \([u:v]=[0:1]\), one must instead use
Eq.~\eqref{eq:eta-polynomial}.  Replacing \(\xi=v/u\) by \(u/v\) without
transforming the vector-field coefficient reverses the chart convention and
can exchange which coefficient detects the point at infinity.
\end{remark}

\section{The exact discriminant--Killing identities}

\subsection{The square of the projective matrix}

Define the binary-quadratic discriminant
\begin{equation}
 \Delta(p):=b^2-4ac.
 \label{eq:discriminant}
\end{equation}
Direct multiplication, with no spectral assumption, gives
\begin{align}
 A_p^2
 &=
 \begin{pmatrix}
 b^2/4-ac&bc/2-bc/2\\
 -ab/2+ab/2&b^2/4-ac
 \end{pmatrix}\\
 &=\boxed{\frac{\Delta(p)}4I_2}.
 \label{eq:A-square}
\end{align}
Equivalently, Cayley--Hamilton for a traceless two-by-two matrix gives
\(A_p^2=-\det(A_p)I_2\), and
\begin{equation}
 \det(A_p)=-\frac{\Delta(p)}4.
 \label{eq:det-A}
\end{equation}

\subsection{Killing normalization from the adjoint trace}

The factor two in AP-E2 is fixed by the canonical Killing form, not by a
rescaling chosen after the fact.

\begin{lemma}[Fundamental-trace formula]
For \(A,B\in\sltwo\),
\begin{equation}
 B_{\rm Kill}(A,B):=\Tr(\ad_A\ad_B)=4\Tr_2(AB).
 \label{eq:killing-trace}
\end{equation}
\end{lemma}

\begin{proof}
Write
\begin{equation}
 A=\begin{pmatrix}h&e\\f&-h\end{pmatrix},
 \qquad
 B=\begin{pmatrix}h'&e'\\f'&-h'\end{pmatrix}.
\end{equation}
Computing the adjoint maps on any fixed \((E,H,F)\) basis gives
\begin{equation}
 \Tr(\ad_A\ad_B)
 =8hh'+4(ef'+fe').
\end{equation}
On the other hand,
\begin{equation}
 4\Tr_2(AB)=4\{2hh'+ef'+fe'\},
\end{equation}
which is identical.  This fixes the normalization of
Eq.~\eqref{eq:killing-trace}.
\end{proof}

Combining Eqs.~\eqref{eq:A-square} and \eqref{eq:killing-trace},
\begin{equation}
 \boxed{
 B_{\rm Kill}(A_p,A_p)
 =4\Tr_2(A_p^2)
 =2\Delta(p).}
 \label{eq:killing-discriminant}
\end{equation}
Writing \(B=\Delta\) would therefore be a normalization error, not a harmless
choice of notation.

\subsection{Second transvectant and polarization}

With the derivative convention used in Route E, the second transvectant of a
quadratic with itself is
\begin{align}
 (p,p)_2
 &:=p_{uu}p_{vv}-2p_{uv}p_{uv}+p_{vv}p_{uu}\\
 &=8ac-2b^2=-2\Delta(p).
 \label{eq:self-transvectant}
\end{align}
Thus
\begin{equation}
 \boxed{(p,p)_2=-B_{\rm Kill}(A_p,A_p).}
 \label{eq:transvectant-killing}
\end{equation}

The bilinear statement is stronger than the diagonal one.  Let
\begin{equation}
 r(u,v)=a'u^2+b'uv+c'v^2
\end{equation}
and polarize the discriminant:
\begin{equation}
 D(p,r):=bb'-2(ac'+a'c),
 \qquad D(p,p)=\Delta(p).
 \label{eq:polar-discriminant}
\end{equation}
Then direct multiplication gives
\begin{equation}
 \boxed{
 B_{\rm Kill}(A_p,A_r)=2D(p,r),
 \qquad (p,r)_2=-2D(p,r)=-B_{\rm Kill}(A_p,A_r).}
 \label{eq:polar-identities}
\end{equation}
This polarization prevents a false proof that happens to reproduce only a
single diagonal example.

\section{Zeros, eigenlines, and the null boundary}

\subsection{Regular and degenerate projective divisors}

Over \(\CC\), a nonzero homogeneous quadratic has a degree-two zero divisor,
counted with multiplicity.  If \(c\ne0\), its finite affine roots are
\begin{equation}
 \xi_{\pm}=\frac{-b\pm\sqrt{\Delta}}{2c}
 \quad(c\ne0),
 \label{eq:affine-roots}
\end{equation}
When \(c=0\) and \(b\ne0\), one root is finite and the other is
\([0:1]\); when \(c=b=0\) and \(a\ne0\), \([0:1]\) is a double root.
The homogeneous formulation handles both cases without a singular formula.

Equation~\eqref{eq:A-square} proves the corresponding Lie-algebra
classification:
\begin{itemize}
\item if \(\Delta\ne0\), the eigenvalues of \(A_p\) are
\(\pm\sqrt{\Delta}/2\), so \(A_p\) is regular semisimple;
\item if \(\Delta=0\) and \(p\ne0\), then \(A_p\ne0\) and \(A_p^2=0\), so
\(A_p\) is rank-one nilpotent;
\item if \(p=0\), then \(A_p=0\), which is neither a nonzero nilpotent orbit
nor a degree-two effective zero divisor.
\end{itemize}

\begin{theorem}[AP-E2 discriminant/contact equivalence]
\label{thm:ap-e2-main}
For a nonzero
\(X_p\in H^0(\CP^1,T_{\CP^1})\), the following are equivalent:
\begin{enumerate}
\item the zero divisor of \(X_p\) is the sum of two distinct projective
points;
\item \(\Delta(p)\ne0\);
\item \(A_p\) is regular semisimple;
\item \(B_{\rm Kill}(A_p,A_p)\ne0\);
\item the Route-E contact norm is nonzero.
\end{enumerate}
The complement inside the nonzero section space is the nilpotent/contact-null
cone.  Each of its points has one double projective zero.  The origin is an
additional zero-section stratum and is excluded from both orbit statements.
\end{theorem}

\begin{proof}
The roots of a nonzero binary quadratic are distinct exactly when
\(\Delta\ne0\).  Equations~\eqref{eq:A-square} and \eqref{eq:det-A} show that
\(\Delta\ne0\) is equivalent to two distinct eigenvalues and hence to regular
semisimplicity.  Equation~\eqref{eq:killing-discriminant} equates this with
nonzero Killing norm.  Section~\ref{sec:route-e-basis} below proves that the
Route-E contact is the same bilinear invariant with fixed nonzero
normalization.  If \(\Delta=0\) and \(p\ne0\), the preceding rank-one
nilpotent argument applies.  A quadratic with zero discriminant is a square
of a nonzero linear form, so its unique projective root has multiplicity two.
Finally, \(p=0\) gives \(A_p=0\) and must be separated before the nonzero
orbit classification is stated.
\end{proof}

\subsection{Finite, zero, and infinite double roots}

The differential of the homogeneous section is
\begin{equation}
 \dd p=(2au+bv)\dd u+(bu+2cv)\dd v.
 \label{eq:first-jet}
\end{equation}
A double projective root is characterized by
\begin{equation}
 p(u,v)=0,
 \qquad 2au+bv=0,
 \qquad bu+2cv=0.
 \label{eq:double-root-jet}
\end{equation}
The following cases explicitly guard both charts.
\begin{center}
\begin{tabular}{@{}lll@{}}
\toprule
Section & Double root & Chart interpretation\\
\midrule
\((2u-v)^2\) & \([1:2]\) & finite nonzero \(\xi=2\)\\
\(v^2\) & \([1:0]\) & finite zero \(\xi=0\)\\
\(u^2\) & \([0:1]\) & \(\eta=0\), the point at infinity\\
\bottomrule
\end{tabular}
\end{center}
All three satisfy Eq.~\eqref{eq:double-root-jet}; restricting to one affine
quadratic formula would miss the last case.

\begin{remark}[Nondegenerate form versus null vector]
The contact matrix is invertible, yet a nonzero \(p\) can obey
\(B(p,p)=0\).  This is ordinary isotropy of a complex symmetric bilinear form,
not degeneracy of the form.  Degeneracy would mean that some nonzero vector
pairs to zero with \emph{every} vector.
\end{remark}

\section{Exact \texorpdfstring{$SL(2)$}{SL(2)} covariance}

Let
\begin{equation}
 g=\begin{pmatrix}\alpha&\beta\\\gamma&\delta\end{pmatrix}
 \in SL(2,\CC),
 \qquad \alpha\delta-\beta\gamma=1,
 \label{eq:g-matrix}
\end{equation}
and define the right substitution
\begin{equation}
 p^g(z):=p(gz).
 \label{eq:right-substitution}
\end{equation}
Expansion gives
\begin{align}
 a^g&=a\alpha^2+b\alpha\gamma+c\gamma^2,\\
 b^g&=2a\alpha\beta+b(\alpha\delta+\beta\gamma)
       +2c\gamma\delta,\\
 c^g&=a\beta^2+b\beta\delta+c\delta^2.
 \label{eq:coefficient-transform}
\end{align}

\begin{proposition}[Conjugacy and invariant discriminant]
Under Eq.~\eqref{eq:right-substitution},
\begin{equation}
 A_{p^g}=g^{-1}A_pg,
 \qquad
 \Delta(p^g)=\Delta(p).
 \label{eq:conjugacy}
\end{equation}
For \(g\in GL(2,\CC)\), the second equality generalizes to
\(\Delta(p^g)=\det(g)^2\Delta(p)\).
\end{proposition}

\begin{proof}
Insert Eq.~\eqref{eq:coefficient-transform} into
Eq.~\eqref{eq:projective-matrix}; matrix multiplication gives the first
identity.  Conjugacy preserves \(\det A_p=-\Delta(p)/4\), proving the
\(SL(2)\) statement.  For a general \(g\), either direct expansion or
factoring \(g\) into a scalar and an \(SL(2)\) matrix gives the determinant
weight two.
\end{proof}

Scaling the section itself by \(\lambda\in\CC\) gives the separate covariance
law
\begin{equation}
 A_{\lambda p}=\lambda A_p,
 \qquad \Delta(\lambda p)=\lambda^2\Delta(p),
 \qquad B(A_{\lambda p},A_{\lambda p})=\lambda^2B(A_p,A_p).
 \label{eq:section-scaling}
\end{equation}
Thus the vanishing locus is projectively well defined even though the
quadratic invariant has weight two.

\subsection{Uniqueness of the invariant contact line}

The contact is not one arbitrary matrix among many.  In the unnormalized
coefficient vector
\begin{equation}
 \bm c=(a,b,c)^T,
 \label{eq:coefficient-vector}
\end{equation}
the infinitesimal \(\mathfrak{sl}_2\) action can be represented by
\begin{equation}
 \rho(H)=\begin{pmatrix}2&0&0\\0&0&0\\0&0&-2\end{pmatrix},
 \quad
 \rho(E)=\begin{pmatrix}0&1&0\\0&0&2\\0&0&0\end{pmatrix},
 \quad
 \rho(F)=\begin{pmatrix}0&0&0\\2&0&0\\0&1&0\end{pmatrix}.
 \label{eq:coefficient-generators}
\end{equation}
Let \(M=M^T\) be a general six-parameter symmetric matrix.  Invariance means
\begin{equation}
 \rho(X)^TM+M\rho(X)=0,
 \qquad X=H,E,F.
 \label{eq:invariance-equations}
\end{equation}
The \(H\) equation permits only \(M_{13}=M_{31}\) and \(M_{22}\).
The \(E\) equation then fixes \(M_{13}=-2M_{22}\); the \(F\) equation adds
no independent freedom.  Hence the solution space is one-dimensional:
\begin{equation}
 M=\kappa
 \begin{pmatrix}
 0&0&-2\\
 0&1&0\\
 -2&0&0
 \end{pmatrix}.
 \label{eq:unique-coefficient-contact}
\end{equation}
At \(\kappa=1\),
\(\bm c^TM\bm c=b^2-4ac=\Delta\).  Equivalently, the exact linear system has
rank five in a six-dimensional space of symmetric matrices.  This is the
first-principles uniqueness of the invariant symmetric contact line.

\section{Route-E basis ledger and the factor of two}
\label{sec:route-e-basis}

\subsection{Spherical coordinates}

Route E uses
\begin{equation}
 \bm x=(x_+,x_0,x_-)^T,
 \qquad
 \Ktr=\frac1{\sqrt3}
 \begin{pmatrix}
 0&0&-1\\
 0&1&0\\
 -1&0&0
 \end{pmatrix}.
 \label{eq:route-e-k}
\end{equation}
The precise map to Eq.~\eqref{eq:binary-quadratic} is
\begin{equation}
 a=\frac{x_-}{\sqrt2},
 \qquad b=-x_0,
 \qquad c=\frac{x_+}{\sqrt2},
 \label{eq:spherical-map}
\end{equation}
or
\begin{equation}
 p_{\bm x}(u,v)
 =\frac{x_-}{\sqrt2}u^2-x_0uv+\frac{x_+}{\sqrt2}v^2.
 \label{eq:spherical-polynomial}
\end{equation}
It follows that
\begin{align}
 \Delta(p_{\bm x})
 &=x_0^2-2x_+x_-,\\
 \bm x^T\Ktr\bm x
 &=\frac1{\sqrt3}(x_0^2-2x_+x_-).
 \label{eq:spherical-contact}
\end{align}
Combining this with Eq.~\eqref{eq:killing-discriminant} yields the frozen
Route-E identity
\begin{equation}
 \boxed{
 B_{\rm Kill}(A_p,A_p)
 =2\Delta(p)
 =2\sqrt3\,\bm x^T\Ktr\bm x.}
 \label{eq:route-e-master}
\end{equation}

The matrix is symmetric and invertible:
\begin{equation}
 \Ktr^2=\frac13I_3,
 \qquad \Ktr^{-1}=3\Ktr,
 \qquad \det\Ktr=-\frac1{3\sqrt3}.
 \label{eq:k-properties}
\end{equation}
Nevertheless,
\begin{equation}
 \bm x=(1,1,1/2)^T
 \quad\Longrightarrow\quad
 \bm x^T\Ktr\bm x=0,
 \qquad \Ktr\bm x\ne0,
 \label{eq:isotropic-example}
\end{equation}
which explicitly distinguishes a nonzero isotropic vector from a degenerate
pairing.

\subsection{Normalized polynomial coefficients}

Now expand the same section in the normalized monomial basis
\begin{equation}
 (u^2,\sqrt2uv,v^2),
 \qquad
 p=y_-u^2+y_0\sqrt2uv+y_+v^2.
 \label{eq:normalized-monomials}
\end{equation}
Thus
\begin{equation}
 \bm y=(y_-,y_0,y_+)^T
 =\begin{pmatrix}a\\b/\sqrt2\\c\end{pmatrix}
 =\frac1{\sqrt2}
 \underbrace{\begin{pmatrix}0&0&1\\0&-1&0\\1&0&0\end{pmatrix}}_{R}
 \bm x.
 \label{eq:y-from-x}
\end{equation}
Because \(R^T\Ktr R=\Ktr\),
\begin{equation}
 \bm y^T\Ktr\bm y=\frac12\bm x^T\Ktr\bm x.
 \label{eq:factor-two-basis}
\end{equation}
Therefore the same Killing invariant has two correct coefficient formulas:
\begin{equation}
 \boxed{
 B=2\sqrt3\,\bm x^T\Ktr\bm x
  =4\sqrt3\,\bm y^T\Ktr\bm y.}
 \label{eq:two-basis-master}
\end{equation}

\begin{center}
\begin{tabularx}{\textwidth}{@{}>{\raggedright\arraybackslash}p{0.23\textwidth}
                                 >{\raggedright\arraybackslash}X
                                 >{\raggedright\arraybackslash}p{0.25\textwidth}@{}}
\toprule
Coordinates & Section convention & Killing formula\\
\midrule
\(\bm c=(a,b,c)\) & \(au^2+buv+cv^2\) &
\(B=2(b^2-4ac)\)\\
\(\bm y=(a,b/\sqrt2,c)\) &
\(y_-u^2+y_0\sqrt2uv+y_+v^2\) &
\(B=4\sqrt3\,\bm y^T\Ktr\bm y\)\\
\(\bm x=(x_+,x_0,x_-)\) &
Eq.~\eqref{eq:spherical-polynomial} &
\(B=2\sqrt3\,\bm x^T\Ktr\bm x\)\\
\bottomrule
\end{tabularx}
\end{center}
Using the spherical coefficient \(2\sqrt3\) on \(\bm y\) loses a factor of
two.  Omitting the \(\sqrt2\) in Eq.~\eqref{eq:spherical-map} changes
\(\Delta=x_0^2-2x_+x_-\) into the wrong expression
\(x_0^2-4x_+x_-\).  Both errors are explicit negative controls in the
regression suite.

\subsection{Complex symmetric is not Hermitian}

The Route-E contact is
\begin{equation}
 \Ktr(\bm x,\bm y)=\bm x^T\Ktr\bm y,
 \label{eq:complex-bilinear}
\end{equation}
which is complex bilinear and symmetric.  Replacing it by
\(\bm x^\dagger\Ktr\bm y\) introduces complex conjugation and destroys the
holomorphic invariant polynomial.  In particular,
\begin{equation}
 (\lambda\bm x)^T\Ktr(\lambda\bm x)
 =\lambda^2\bm x^T\Ktr\bm x,
\end{equation}
whereas the Hermitian expression scales as \(|\lambda|^2\).  These are
different representations and cannot be exchanged on complex samples.

\section{The H3 versus H3\texorpdfstring{$^+$}{+} boundary}

An invariant nondegenerate symmetric bilinear form does not by itself imply a
nonzero Killing form.  Let \(\mathfrak a=\CC e\) be a one-dimensional abelian
Lie algebra.  Then
\begin{equation}
 [e,e]=0,
 \qquad \ad_e=0,
 \qquad B_{\rm Kill}(e,e)=\Tr(\ad_e^2)=0.
 \label{eq:abelian-killing}
\end{equation}
However, the form \(C(e,e)=1\) is symmetric, invariant, and nondegenerate,
because every bracket vanishes.  Hence the original H3 condition ``there is a
nondegenerate invariant symmetric pairing'' admits this genus-one/abelian
counterexample.  Exact selection of the three-dimensional
\(\sltwo\) carrier requires the stronger H3$^+$ condition that the contact is
the nondegenerate adjoint-trace/Killing line.

\begin{corollary}[Precise AP-E2 scope]
Theorem~\ref{thm:ap-e2-main} proves an exact equivalence on the H3$^+$-selected
\(\CP^1\) branch.  It neither derives H3$^+$ from H3 nor rules out the abelian
H3 branch.
\end{corollary}

\section{The 30-check fail-closed regression}

\subsection{Reproducibility contract}

The executable card is
\begin{center}
\path{route_f/code/verify_ap_e2_discriminant_regression.py}.
\end{center}
From the repository root, run
\begin{center}
\texttt{python3 route\_f/code/verify\_ap\_e2\_discriminant\_regression.py}.
\end{center}
It writes
\begin{center}
\path{route_f/output/ap_e2_discriminant_regression.json}\\
\path{route_f/output/ap_e2_discriminant_regression.md}.
\end{center}
The seed is \(20260714\).  Exact tests use
\path{fractions.Fraction}; irrational and complex tests use
\path{mpmath 1.4.1} at 100 decimal places.  The program writes its complete
report before returning exit status one if any check fails.

\subsection{Coverage by group}

\begin{longtable}{@{}>{\raggedright\arraybackslash}p{0.12\textwidth}
                       >{\centering\arraybackslash}p{0.10\textwidth}
                       >{\raggedright\arraybackslash}p{0.68\textwidth}@{}}
\toprule
Group & Checks & Fail-closed content\\
\midrule
\endfirsthead
\toprule
Group & Checks & Fail-closed content\\
\midrule
\endhead
A0 & 4 & Critical source hashes, invariant-card stable digest, Audit-0 builder
binding, and frozen theorem/proof-block hash.\\
A1 & 10 & Exact rational \(A^2\), Killing, transvectant, polarized contact,
one-dimensional invariant line, \(\det(z,Az)=p(z)\), \(SL(2,\mathbb Z)\)
invariance and conjugacy, and section-scaling weight.\\
A2 & 5 & Finite and infinite double roots, real/complex/asymmetric regular
roots, eigenline agreement, vanishing first jet, zero-section guard, and a
\(10^{-70}\) near-null perturbation.\\
A3 & 4 & Route-E \(\Ktr\) normalization, inverse and isotropic-vector guard,
spherical/normalized-basis factor two, and full \(H,E,F\) invariance.\\
A4 & 1 & 256 seeded complex samples evaluated at 100 decimal places.\\
A5 & 4 & Required rejection of missing \(\sqrt2\), missing Killing factor two,
wrong normalized-basis coefficient, and Hermitian substitution.\\
A6 & 2 & The abelian H3 counterexample and the immutable non-promotion Boolean.\\
\midrule
Total & 30 & Every check must pass; otherwise the process exits nonzero.\\
\bottomrule
\end{longtable}

The exact randomized layer uses 256 seeded triples in \(\QQ^3\), 256 seeded
pairs for the polarized identity, 256 exact projective fixed-point tests, and
256 determinant-one integer substitutions.  Since all operations are in
\(\QQ\), ``zero residual'' means an exact equality, not a floating tolerance.

The complex layer separately samples 256 points in \(\CC^3\).  Its largest
intrinsic normalized residual is
\begin{equation}
 1.67930436345868921\times10^{-101}.
 \label{eq:intrinsic-residual}
\end{equation}
The Route-E JSON card stores ordinary double-precision decimals, so its
separate maximum source-card residual is
\begin{equation}
 6.14719575160384581\times10^{-17},
 \label{eq:source-residual}
\end{equation}
well below the declared \(5\times10^{-15}\) source tolerance.  Keeping these
two residuals separate prevents the source serialization precision from
masquerading as a limitation of the 100-digit identity test.

For the near-null probe
\begin{equation}
 p_\varepsilon=u^2+2uv+(1+\varepsilon)v^2,
 \qquad \varepsilon=10^{-70},
\end{equation}
the audit resolves
\begin{equation}
 \Delta(p_\varepsilon)=-4\times10^{-70}
\end{equation}
and obtains zero recorded matrix-identity residual at the working precision.
This tests the boundary from the regular side instead of evaluating only the
exact cone.

\subsection{Negative controls}

A regression is incomplete if it tests only the intended formula.  AP-E2
therefore demands that four intentionally wrong conventions disagree on a
fixed generic complex point.
\begin{center}
\begin{tabularx}{\textwidth}{@{}>{\raggedright\arraybackslash}X
                                 >{\raggedright\arraybackslash}p{0.27\textwidth}
                                 >{\raggedleft\arraybackslash}p{0.16\textwidth}@{}}
\toprule
Wrong convention & Expected diagnosis & Residual\\
\midrule
\(a=x_-,\ b=-x_0,\ c=x_+\) & missing \(\sqrt2\) & \(1.0\times10^1\)\\
\(B=\Delta\) & missing Killing factor two & \(1.2\times10^1\)\\
\(B=2\sqrt3\,\bm y^T\Ktr\bm y\) & spherical factor used on normalized
coefficients & \(1.2\times10^1\)\\
\(\bm x^\dagger\Ktr\bm x\) & Hermitian form substituted for bilinear contact
& \(1.56205\times10^1\)\\
\bottomrule
\end{tabularx}
\end{center}
Detection of each wrong formula is a pass.  Accidental agreement is a failure.
An additional injected bad theorem hash was tested during development; it
produced 29/30 and a nonzero process exit before the valid artifact was
regenerated.

\subsection{Current result and provenance lock}

The current card reports
\begin{equation}
 \boxed{30/30\ \text{checks pass}.}
 \label{eq:thirty-pass}
\end{equation}
Its non-promotion field remains
\begin{center}
\nopromote.
\end{center}
The frozen theorem/proof block has SHA-256
\begin{center}
\path{3999b88f3fe7e396948c9695939ec5826d74cfc6d06f6b26ed05d3b0e0c4d094}.
\end{center}
The audit also recomputes the stable digest of the Route-E invariant card and
checks that the card-recorded Audit-0 builder hash equals the current builder.
Thus a green result requires both correct arithmetic and an intact convention
source chain.

\section{Claim ledger and handoff to AP-E3}

\begin{longtable}{>{\raggedright\arraybackslash}p{0.42\textwidth}
                  >{\raggedright\arraybackslash}p{0.17\textwidth}
                  >{\raggedright\arraybackslash}p{0.32\textwidth}}
\toprule
Claim & AP-E2 status & Evidence or missing input\\
\midrule
\endfirsthead
\toprule
Claim & AP-E2 status & Evidence or missing input\\
\midrule
\endhead
Global vector fields on \(\CP^1\) are binary quadratics &
\status{proved} & Two-chart regularity, Eqs.~\eqref{eq:eta-transform}--
\eqref{eq:q-quadratic}.\\
Zeros of \(p\) equal eigenlines of \(A_p\) &
\status{proved} & Determinant identity \eqref{eq:wedge-identity}.\\
Two distinct zeros, regular semisimplicity, and non-null contact are
equivalent & \status{proved on carrier} &
Theorem~\ref{thm:ap-e2-main}.\\
The null cone consists of nonzero double-zero/nilpotent sections &
\status{proved} & Matrix square and first-jet tests; zero section excluded.\\
The Route-E factor is \(B=2\sqrt3\,x^T\Ktr x\) &
\status{proved} & Killing trace plus spherical map.\\
Normalized coefficients use \(B=4\sqrt3\,y^T\Ktr y\) &
\status{proved} & Exact basis map \eqref{eq:y-from-x}.\\
H3 alone selects the three-dimensional carrier &
\status{refuted} & One-dimensional abelian invariant-pairing counterexample.\\
H3$^+$ has a microscopic dynamical origin &
\status{open} & Requires a correlator, messenger, anomaly, or UV derivation.\\
The two divisor points are physical particles, signs of energy, or families &
\status{not derived} & No action or observable map supplies this identification.\\
The projective sector has level magnitude \(|k|=2\) &
\status{candidate magnitude derived; sign/UV open} & The AP-E3 Mott--Hund
pair gives a ket line \(\cO(-2)\) and dual prequantum line \(\cO(2)\);
the signed orientation and exactly-two UV rule remain absent.\\
\bottomrule
\end{longtable}

AP-E2 is therefore mathematically closed at the regression level.  AP-E3
should not repeat this discriminant calculation.  It should instead test
several genuinely distinct microscopic routes to level two:
\begin{enumerate}
\item derive a charge-two collective coordinate and its Berry one-form from a
fully normalized microscopic action;
\item integrate out explicitly specified gapped fields and compute the induced
worldline Chern--Simons/Berry level, including anomaly and sign checks;
\item derive a pair-selection rule that leaves a level-two composite sector
while keeping the macroscopic Q-ball charge separate;
\item reject every proposal for which the induced level is not integral, is
shifted by an uncancelled anomaly, or has a Landau-level gap comparable to the
portal or background scale.
\end{enumerate}
None of these AP-E3 mechanisms can be inferred from
\(\Delta=B/2\) alone.  That separation is the principal scientific output of
AP-E2.

\section{Conclusion}

Starting only from the two charts of \(\CP^1\), AP-E2 derives the complete
chain
\begin{equation}
 \det(z,A_pz)=p(z),
 \qquad
 A_p^2=\frac\Delta4I,
 \qquad
 B_{\rm Kill}=2\Delta,
 \qquad
 (p,p)_2=-2\Delta.
\end{equation}
It proves that the discriminant, projective two-center divisor,
regular/nilpotent stratification, Killing norm, and Route-E contact are one
invariant statement on the conditional \(\CP^1\) carrier.  The chart at
infinity, zero-section exception, complex-bilinear character, and
spherical/normalized-basis factor two are all explicit rather than implicit.

The 30-check regression now freezes these conclusions against known
convention failures and source drift.  This completes AP-E2 without promoting
the bridge.  AP-E3 now supplies a conditional Mott--Hund pair with
\(|k|=2\); its displayed aligned action has \(k=-2\), while the desired sign,
the physical origin of H3$^+$, and the UV principle selecting exactly two
constituents remain logically separate open problems.

\end{document}
